\begin{table}[htb]
\centering
\small
\setlength{\tabcolsep}{4pt}
\caption[Espectro de Fourier medido]{Espectro de Fourier medido por FFT 2D contra o limite da Tabela~\ref{tab:espectros}, com $d = 2$ atributos.}
\label{tab:espectro}
\begin{tabular}{lccccccc}
\toprule
Codificação & \shortstack{$\omega_{\max}$\\previsto} & \shortstack{$\omega_{\max}$\\em $x_1$} & \shortstack{$\omega_{\max}$\\em $x_2$} & \shortstack{Termos\\previstos} & \shortstack{Termos\\medidos} & Cruzados & $p$ \\
\midrule
\textit{Angle} & 1 & 1 & 1 & 9 & 6 & 4 & 12 \\
\textit{Re-uploading} ($R = 1$) & 1 & 0 & 1 & 9 & 2 & 0 & 6 \\
\textit{Re-uploading} ($R = 2$) & 2 & 1 & 2 & 25 & 10 & 8 & 12 \\
\textit{Re-uploading} ($R = 3$) & 3 & 2 & 3 & 49 & 35 & 24 & 18 \\
\bottomrule
\end{tabular}
\\[2pt]
{\footnotesize A teoria dá um limite superior: nenhuma energia aparece fora de $\{-R, \dots, R\}^d$, mas nem todos os termos permitidos aparecem. \textit{Termos medidos}: pares $(\omega_1, \omega_2)$ com energia na FFT 2D; \textit{Cruzados}: os que dependem dos dois atributos. Codificações \textit{amplitude} e ZZ não aparecem porque a saída não é uma série de Fourier nos atributos. Em $x_1$, a frequência máxima é $R - 1$: a última inserção de $x_1$ atua só no qubit 0 e não afeta a leitura de $Z_0$.}
\\[2pt]
{\footnotesize Fonte: Elaborada pelo autor.}
\end{table}
